\section{A Game of Robots}
\label{sec:g-robots}
\source{ICPC 2018--19 Online Preliminary (PYQ, ROBOGAME)}{Medium}

\problemstatement{A 1-D grid of $n\le50$ cells is given as a string: \texttt{.}
is empty, a digit $x$ is a robot that moves back and forth inside the cells
$[i-x,i+x]$ (clipped to the grid). You choose each robot's initial direction;
the robots choose their own start times and \emph{want} to collide. Output
\texttt{safe} if you can prevent every collision, else \texttt{unsafe}.}

\begin{patternbox}
The robots control their start times, so direction choices cannot stop two
robots whose ranges share a cell. Reduce to \textbf{interval overlap}, then a
difference array.
\end{patternbox}

\subsection*{Approach and intuition}

If the ranges of two robots intersect in some cell $c$, each robot passes
through $c$ periodically; by delaying one robot's start by the right amount
they are at $c$ at the same moment, whatever directions we chose. If no two
ranges intersect, the robots never share a cell at all. So:
\[
  \text{safe}\iff\text{no cell is covered by two ranges.}
\]
A robot with $x=0$ occupies just its own cell --- it still counts.

\subsection*{Algorithm}
\begin{algorithmic}[1]
\State $diff[\cdot]\gets0$
\For{each robot at $i$ with range $x$} \State $diff[\max(0,i-x)]$\texttt{++}; $diff[\min(n-1,i+x)+1]$\texttt{--} \EndFor
\State \Return \texttt{unsafe} if some prefix sum $\ge2$, else \texttt{safe}
\end{algorithmic}

\dryrun{\texttt{.2...2..} (0-indexed): ranges $[0,3]$ and $[3,7]$ share cell
$3$ $\Rightarrow$ unsafe. \texttt{.2.....}: one robot $\Rightarrow$ safe.
\texttt{1.1.1.}: ranges $[0,1]$, $[1,3]$ overlap $\Rightarrow$ unsafe.
\checkmark}

\subsection*{C++17 implementation}
\cppfile{greedy/24_game_of_robots.cpp}

\begin{complexitybox}
\textbf{Time:} $O(n)$ per test. \textbf{Space:} $O(n)$.
\end{complexitybox}

\begin{pitfallbox}
\begin{itemize}
  \item Clip ranges to the grid \emph{before} checking overlap.
  \item The long story about directions is a red herring; the key sentence
        is ``they can coordinate \dots\ and decide when to start''.
\end{itemize}
\end{pitfallbox}
